% Part: first-order-logic
% Chapter: axiomatic-deduction
% Section: axioms-rules-quantifiers

\documentclass[../../../include/open-logic-section]{subfiles}

\begin{document}

\olfileid{fol}{axd}{qua}

\olsection{Aksioma lan Aturan kanggo Kuantor}

\begin{defn}[Aksioma kanggo kuantor]
\emph{Aksioma} kang ngatur kuantor yaiku kabeh conto saka wujud iki:
\begin{align}
\ollabel{ax:q1} & \lforall[x][!B] \lif !B(t), \\
\ollabel{ax:q2} & !B(t) \lif \lexists[x][!B].
\end{align}
kanggo saben term tertutup~$t$.
\end{defn}

\begin{defn}[Aturan kanggo kuantor]
\item Yen $!B \lif !A(a)$ wis muncul ing !!{derivation} lan $a$ ora
  muncul ing~$\Gamma$ utawa~$!B$, mula $!B \lif
  \lforall[x][!A(x)]$ dadi langkah inferensi kang bener.
\item Yen $!A(a) \lif !B$ wis muncul ing !!{derivation} lan $a$ ora
  muncul ing~$\Gamma$ utawa~$!B$, mula $\lexists[x][!A(x)] \lif
  !B$ dadi langkah inferensi kang bener.
\end{defn}

Kita bakal nyingkat salah siji saka aturan iki dadi ``\QR.''

\end{document}
